\addtocontents{toc}{\protect\newpage}
\chapter{Профілі безпеки NIST і їх оцінка відповідності}

Комплексна система захисту інформації (КСЗІ) є невід'ємним елементом побудови
державних інформаційних систем України. Для отримання \emph{Атестата відповідності}
система повинна пройти чітко регламентовану процедуру, яка охоплює розробку цільового
профілю безпеки, імплементацію технічних і криптографічних засобів захисту та
незалежну державну експертизу. Нормативно-правовою базою цього процесу слугують
нові стандарти ТЗІ, повністю гармонізовані з міжнародним каталогом \textbf{NIST SP~800-53}.

\section{Безпекові переваги платформи Erlang}

Платформа \texttt{synrc/ca} розроблена на базі віртуальної машини \textbf{BEAM (Erlang/OTP)},
створеної спеціально для безперервних телекомунікаційних систем. Це архітектурне рішення
виключає цілі класи критичних вразливостей (CVE) та дозволяє забезпечити нативну підтримку
ключових контролів безпеки NIST, які принципово недосяжні на віртуальних машинах CLR та JVM:

\begin{itemize}
    \item \textbf{SI-16 — Захист пам'яті.} Відсутність ручного керування пам'яттю повністю
          виключає вразливості типу \emph{Buffer~Overflow} та \emph{Use-After-Free}, які є
          найчастішим вектором атак у класичних C/C++\, серверах.

    \item \textbf{SC-39 — Ізоляція процесу.} Мільйони легковагових акторів (процесів) ізольовані
          на рівні віртуальної машини. Збій або компрометація одного процесу ніколи не вплине на
          інші, забезпечуючи абсолютну стійкість. Кожен BEAM-процес має власний збирач сміття,
          що відрізняє цю модель від JVM та CLR.

    \item \textbf{SI-13 — Запобігання збоїв.} Філософія <<Let it crash>> та потужні дерева
          супервізорів (Supervision Trees) забезпечують автоматичне відновлення стану системи
          за мілісекунди без простоїв.

    \item \textbf{AC-25 — Опорний монітор.} Маршрути, правила ABAC та критичні реєстри
          компілюються прямо в байт-код (memory-resident). Це виключає затримки бази даних і
          гарантує неможливість обходу політик (non-bypassable).

    \item \textbf{SC-5 — Захист від відмови в обслуговуванні.} На відміну від JVM/CLR із
          паузами Stop-The-World, кожен BEAM-процес має мікро-GC. Витискальна багатозадачність
          гарантує Soft~Real-time відгук і захист від Application DoS-атак.

    \item \textbf{SI-2 — Усунення вразливостей (Hot Code Swapping).} Дозволяє замінювати
          скомпільований код (патчі) <<на льоту>> без рестарту VM. Вразливості закриваються або
          логіка змінюється без розриву активних з'єднань клієнтів.
\end{itemize}

\section{Гомотопічне кодування ASN.1 DER}

Стандарт \textbf{ASN.1 DER} обраний як фундаментальний формат серіалізації для всіх
API-ендпойнтів екосистеми \texttt{synrc/ca}. На відміну від Protobuf / gRPC та JSON, які
є вразливими до маніпуляцій через множинність варіантів кодування одного й того ж об'єкта,
гомотопічний (канонічний) \textbf{DER} математично гарантує \emph{єдине унікальне байтове
представлення} для кожної структури даних. Це повністю виключає атаки типу
\emph{Parser~Differentials} (наприклад, Bleichenbacher's forgery) та забезпечує виконання
наступних вимог безпеки:

\begin{itemize}
    \item \textbf{SI-10 — Перевірка вводу інформації.}
          Сувора перевірка вхідних даних відкидає будь-які нестандартні (BER) або надлишкові
          кодування, унеможливлюючи вразливості до маніпуляцій (malleability).

    \item \textbf{SI-7 — Цілісність інформації.}
          Цифрові підписи та хешування залишаються валідними лише за умови бієктивної
          (гомотопічної) серіалізації, захищаючи від прихованої модифікації байтового представлення.

    \item \textbf{SC-13 — Криптографічний захист.}
          Вимоги стандартів (X.509, CMS, PKCS) для безпечного застосування криптографії та
          збереження довіри до інфраструктури відкритих ключів (PKI).
\end{itemize}

\section{Профілі безпеки КСЗІ}

Система \texttt{synrc/ca} підтримує автоматичну побудову профілів Комплексної системи захисту
інформації (КСЗІ) на основі нормативно-правових актів у галузі технічного захисту інформації
(ТЗІ НПА). Це дозволяє автоматизувати створення документації та мапінг вимог законодавства на
технічні налаштування. Усі сертифікати видаються з імпринтами OID відповідних профілів безпеки.

\subsection{Імплементація опорного монітора (AC-25)}

Для гарантування цілісності та неможливості обходу (non-bypassable), динамічні політики безпеки
(ABAC) не зберігаються у традиційних базах даних. Натомість в ERP/1 Сертифікати матриця доступу
та політики компілюються безпосередньо у незмінний (immutable) BEAM байт-код віртуальної машини
Erlang, функціонуючи як memory-resident механізм. Це архітектурне рішення забезпечує виконання
вимог NIST щодо захисту політик від модифікації (tamper-proof) та повністю нівелює класичні
вектори атак на корупцію пам'яті, притаманні традиційним ANSI C99 реалізаціям, усуваючи
відповідні класи вразливостей (CVE).

\subsection{Стандарти NIST та ISO}

Застосування контролів базується на наступних міжнародних документах:

\begin{itemize}
    \item \textbf{NIST SP 800-53} --- Контроли безпеки та приватності для федеральних інформаційних систем і організацій.
    \item \textbf{NIST SP 800-53A} --- Оцінювання засобів контролю безпеки в організаціях.
    \item \textbf{NIST SP 800-37} --- Структура управління ризиками (Risk Management Framework, RMF).
    \item \textbf{FIPS 199 / FIPS 200} --- Категоризація та мінімальні вимоги до безпеки федеральних систем.
    \item \textbf{NIST SP 800-162} --- Настанови щодо систем управління доступом на основі атрибутів (ABAC).
    \item \textbf{ISO/IEC 27001} --- Системи управління інформаційною безпекою~--- вимоги.
    \item \textbf{ISO/IEC 27002} --- Інформаційні технології~--- методи захисту~--- контролі інформаційної безпеки.
\end{itemize}

Параметри безпеки включають: автоматизовану прив'язку політик безпеки до метрик;
динамічний контроль доступу на основі атрибутів; неперервний аудит, логування подій та моніторинг;
безперервну верифікацію та валідацію криптографічних засобів і сертифікатів.

\section{Походження архітектури контролів: NIST vs ISO}

Класифікація сімейств заходів захисту, яка лежить в основі платформи (абревіатури AC, AU, SC,
IA тощо), \emph{не є частиною} міжнародних стандартів ISO/IEC~27001 або 27002. Ці стандарти мають
інший підхід до групування. Архітектура \texttt{synrc/ca} базується на еталонному каталозі
\textbf{NIST SP 800-53} (США). Саме ця гранулярна база контролів була нещодавно повністю
адаптована Державною службою спеціального зв'язку та захисту інформації України у новому
національному стандарті \textbf{НД ТЗІ 3.6-006-24}. Шлях NIST обраний тому, що він забезпечує
максимальну деталізацію, необхідну для побудови military-grade систем та атестації державних реєстрів.

\subsection{Державний сектор (НД ТЗІ 3.6-006-24)}

У державному секторі України захист інформації регламентується стандартом НД ТЗІ 3.6-006-24,
повністю гармонізованим з NIST SP 800-53. Процедура побудови КСЗІ передбачає:

\begin{enumerate}
    \item Розробку та затвердження \textbf{Цільового профілю безпеки} системи на основі Базового
          або Галузевого профілю.
    \item Впровадження засобів технічного і криптографічного захисту, що відповідають обраним
          заходам контролю.
    \item Проведення державної експертизи КСЗІ (або декларування відповідності) для отримання
          \textbf{Атестата відповідності} від Держспецзв'язку.
\end{enumerate}

Нормативна документація: НД ТЗІ 3.6-006-24, НД ТЗІ 2.3-025-24 (Том 1, Том 2, Том 3).

\subsection{Приватний сектор (NIST SP 800-53)}

Для корпоративного сектору, фінтеху та критичної інфраструктури еталоном є пряме застосування
сімейства стандартів NIST. Сертифікація базується на Risk Management Framework (NIST SP 800-37)
і включає:

\begin{enumerate}
    \item Категоризацію системи за FIPS 199/200 та вибір базового профілю контролів (Baseline)
          за NIST SP 800-53.
    \item Розробку System Security Plan (SSP) та імплементацію архітектури безпеки і засобів
          моніторингу.
    \item Незалежну оцінку контролів за NIST SP 800-53A та формальне надання внутрішнього дозволу
          на експлуатацію (\textbf{Authority to Operate~--- ATO}).
\end{enumerate}

\section{Процес управління ризиками (Risk Management Framework)}

Життєвий цикл побудови та підтримки комплексної системи безпеки за методологією NIST (зокрема згідно з \textbf{NIST SP 800-37})
визначається структурою управління ризиками (Risk Management Framework, RMF). Цей безперервний процес гарантує,
що заходи безпеки відповідають архітектурі, бізнес-вимогам та поточним загрозам, і складається з 6 базових кроків:

\begin{enumerate}
    \item \textbf{Категоризація інформаційної системи (Categorize Information System).} На цьому етапі визначається рівень критичності системи (Low, Moderate, High) на основі аналізу наслідків від потенційної втрати конфіденційності, цілісності чи доступності інформації (відповідно до FIPS 199 та FIPS 200).
    \item \textbf{Вибір заходів безпеки (Select Security Controls).} На основі визначеної категорії обирається базовий набір контролів безпеки (Baseline) з каталогу NIST SP 800-53, який потім адаптується (tailoring) під специфічні умови експлуатації та ризики організації.
    \item \textbf{Впровадження заходів безпеки (Implement Security Controls).} Технічна та організаційна імплементація обраних контролів у системі. Розробляється детальний System Security Plan (SSP), що описує, як саме реалізовано кожен конкретний захід.
    \item \textbf{Оцінювання заходів безпеки (Assess Security Controls).} Незалежна оцінка та аудит впроваджених контролів за допомогою методології NIST SP 800-53A. Визначається, наскільки коректно імплементовані захисні механізми та чи досягають вони бажаного результату.
    \item \textbf{Авторизація інформаційної системи (Authorize Information System).} Офіційне прийняття залишкових ризиків керівництвом організації (Authorizing Official) на основі результатів оцінювання та надання формального дозволу на експлуатацію системи (Authority to Operate, ATO).
    \item \textbf{Моніторинг заходів безпеки (Monitor Security Controls).} Безперервне відстеження ефективності контролів, фіксація змін у системі та її середовищі, аналіз інцидентів безпеки та періодичне оновлення SSP. Цей крок гарантує, що система зберігає відповідність вимогам протягом усього свого життєвого циклу.
\end{enumerate}

\section{Принцип розділення обов'язків}

Відповідно до нормативної бази України (НД ТЗІ 2.6-001-11) та найкращих світових практик
(принцип \emph{Maker/Checker}), побудова та сертифікація КСЗІ є строго \textbf{двоетапною
процедурою}. Вона вимагає залучення двох абсолютно незалежних провайдерів-ліцензіатів для
унеможливлення конфлікту інтересів.

\subsection{Етап 1 --- Розробка профілю безпеки (Провайдер №1)}

Провайдер №1 проводить оцінку ризиків, обстеження середовища та створює технічні інструкції~---
\textbf{Технічне завдання (Цільовий профіль безпеки)}. Здійснює фізичну імплементацію архітектури
безпеки: налаштовує криптографію, Reference Monitor (AC-25), встановлює сертифікати та здає
систему в дослідну експлуатацію.

\subsection{Етап 2 --- Атестація відповідності (Провайдер №2)}

Провайдер №2 отримує Технічне завдання від розробника та розробляє \textbf{Програму експертизи}.
Проводить незалежний технічний аудит (атестація відповідності), інструментальне тестування та перевірку імплементованих
контролів наживо. У разі успіху видає Експертний висновок для Держспецзв'язку.

\section{Каталог профілів безпеки ERP/1}

Нижче наведено перелік розроблених профілів безпеки, кожен з яких визначає цільовий або
галузевий набір контрольних елементів на базі НД ТЗІ 2.3-025-24 та НД ТЗІ 3.6-006-24:

\begin{itemize}
    \item \textbf{FULL DIRECTORY (1123 контроли)}~--- повний каталог контрольних елементів
          із трьохтомника НД ТЗІ 2.3-025-24. Еталонний документ для побудови будь-якого
          галузевого або цільового профілю.

    \item \textbf{OPEN \& CONFIDENTIAL --- 126 контролів}~--- базовий профіль безпеки відкритої
          та конфіденційної інформації, затверджений Наказом Адміністрації Держспецзв'язку №409
          від 30.06.2025.

    \item \textbf{OFFICIAL --- 155 контролів}~--- базовий профіль безпеки службової інформації,
          затверджений Наказом Адміністрації Держспецзв'язку №419 від 02.07.2025.

    \item \textbf{COURT INDUSTRY PROFILE (190 контролів)}~--- галузевий профіль для судової
          системи. Описує розширений набір базових контролів для обробки чутливої інформації в
          державних реєстрах та судах, враховуючи вимоги високої доступності та цілісності даних
          на базі НД ТЗІ 3.6-006-24.

    \item \textbf{DIGITAL UKRAINIAN GOVERNMENT INDUSTRY PROFILE (350 контролів)}~---
          галузевий профіль безпеки Міністерства цифрової трансформації України. Описує
          розширений набір базових контролів для обробки чутливої інформації в державних реєстрах
          та судах на базі НД ТЗІ 3.6-006-24.

    \item \textbf{ERP/1 SECURITY TARGET}~--- цільовий профіль для телекомунікаційної системи
          ERP/1. Описує імплементацію управління доступом на основі атрибутів (ABAC), управління
          інформаційними потоками (BPE), інфраструктури відкритих ключів (PKI) та надійного
          аудиту (KVS) для безпеки критичних бізнес-процесів організації.

    \item \textbf{X.509 CMS MESSAGING SECURITY TARGET}~--- цільовий профіль для систем
          корпоративного месенджера. Базується на криптографічних протоколах X.509 та CMS
          (Cryptographic Message Syntax), забезпечуючи наскрізне шифрування (E2EE), строгу
          автентифікацію та гарантію невідмовності авторства для юридично значущих комунікацій.

    \item \textbf{VPN SECURITY TARGET}~--- цільовий профіль для інфраструктури віртуальних
          приватних мереж (VPN) та центрів сертифікації (PKI). Визначає механізми тунелювання,
          автентифікації вузлів (SC-8) і захисту мережевого трафіку на транспортному рівні для
          захисту каналів зв'язку від перехоплення та модифікації.
\end{itemize}

\newpage
\section{Архітектура EUDI та децентралізована PKI}

EUDI~--- це децентралізований PKIX з контролем атрибутів на рівні ABAC, що використовує JSON
для кодування та HTTP як транспорт. Система визначає наступних учасників:

\begin{itemize}
    \item \textbf{eIDAS Node}~--- державний центр сертифікації (CA);
    \item \textbf{EUDI Verifier}~--- перевірка Verifiable Presentations;
    \item \textbf{EUDI Wallet (Holder)}~--- iOS/Android застосунок для зберігання та презентації
          облікових даних;
    \item \textbf{EUDI Provider (Issuer)}~--- OpenID для Verifiable Credentials (PID, QEAA);
    \item \textbf{Personal Identification Data (PID) Provider}~--- державне підприємство <<Дія>>;
    \item \textbf{QEAA / EAA}~--- кваліфіковані та некваліфіковані атестати атрибутів;
    \item \textbf{Qualified Electronic Signature Provider (QP)}~--- кваліфіковані сертифікати (QC).
\end{itemize}

\subsection{Holder, Issuer, Verifier}

В екосистемі OpenID4VC Verifier та Issuer зв'язані опосередковано через життєвий цикл облікових
даних, а взаємодія між ними в основному здійснюється через Holder. Ця архітектура забезпечує
довіру без прямих, постійних відносин між Verifier та Issuer, дотримуючись принципів
конфіденційності та децентралізації.

Верифікатор не контактує безпосередньо з видавцем під час типових операцій, якщо тільки не
потрібна перевірка статусу. Holder виступає посередником, забезпечуючи конфіденційність та
контроль над даними, що передаються. EUDI Wallet виступає як Holder, а QEAA, EAA, PIP (TSPs)~---
як EUDI Providers або Issuers. EUDI Verifier виконує перевірку статусу облікових даних та
верифікацію презентацій.

\subsection{PKIX vs OpenID4VC}

Модель EUDI має схожість із PKIX: так само, як особа використовує підписаний набір атрибутів
(сертифікат X.509 зі CSR-атрибутами) для автентифікації та авторизації в PKI, провайдер
OpenID4VC (PIP) охоплює набір атрибутів (цифрова презентація claims) та видає електронні
документи у форматі mDOC для EUDI Wallet.

Однак, на відміну від PKIX з його централізованою моделлю, EUDI забезпечує розподілену модель
без єдиного кореневого CA, де всі сторони пов'язані криптографічно. EUDI також здійснює більш
тонкий та строгий контроль над атрибутами (claims), подібно до моделі ABAC.

CRL та OCSP можуть створювати проблеми конфіденційності, оскільки передбачають запит до CA, що
потенційно розкриває активність користувача. OpenID4VC пом'якшує це, дозволяючи Holder
керувати процесом, а деякі реалізації уникають перевірки статусу в режимі реального часу,
включаючи криптографічні докази безпосередньо в облікові дані.

\section{Відповідність вимогам GDPR та захист персональних даних}
\index{GDPR}
При інтеграції державних ІКС України з європейським цифровим простором обов'язковим є дотримання вимог регламенту GDPR (General Data Protection Regulation). Платформа ERP/1 реалізує наступні принципи GDPR:
\begin{itemize}
    \item \textbf{Privacy by Design та Privacy by Default}: Обробка персональних даних мінімізована на рівні протоколів. Усі транзакційні повідомлення в Cartulary (Картулярії) деперсоналізовані та використовують хеш-ідентифікатори користувачів.
    \item \textbf{Обмеження цілей та мінімізація даних}: Атрибути користувача в каталогах LDAP зберігаються за моделлю ABAC. Сервіси отримують доступ лише до тих полів даних, які необхідні для виконання конкретної бізнес-функції (наприклад, перевірки віку без передачі дати народження за технологією Zero-Knowledge Proofs).
    \item \textbf{Право на забуття (Right to be Forgotten)}: В архітектурі незмінних реєстрів (на базі блокчейну або захищеного Картулярію) фізичне видалення записів є неможливим через вимоги цілісності ланцюжка. ERP/1 вирішує цю суперечність за допомогою \emph{криптографічного стирання} (Cryptographic Erasure) --- персональні дані шифруються унікальним симетричним ключем суб'єкта, і при реалізації права на забуття цей ключ безповоротно знищується з IAS, що робить розшифрування персональних даних неможливим.
\end{itemize}

\section{Тестові питання та завдання}
\begin{enumerate}
    \item Які особливості віртуальної машини BEAM (Erlang) забезпечують нативне виконання вимог NIST SP 800-53 щодо захисту пам'яті (SI-16) та ізоляції процесів (SC-39)?
    \item Чому гомотопічне (канонічне) кодування DER є більш безпечним для цифрових підписів порівняно з гнучким кодуванням BER чи JSON?
    \item Опишіть ролі Holder, Issuer та Verifier в архітектурі EUDI Wallet та порівняйте їх із класичною моделлю сертифікатів X.509.
    \item Як реалізується вимога GDPR щодо «права на забуття» в системах з незмінними реєстрами (immutable ledgers)?
    \item \textbf{Практичне завдання}: Опишіть перелік OID атрибутів для сертифіката користувача, які необхідні для авторизації за моделлю ABAC у системі доступу до медичних карток.
\end{enumerate}
